\section{Current Default Formulas}

This section documents the pre-refactor serving formulas that still remain
available through compatibility shims. The relevant implementations live in
\texttt{calculate\_scaled\_serving\_duration(...)} in
\texttt{astra-sim/workload/ServingUtils.cc} and in the legacy-compatible paths
of \texttt{ServingCostModel::estimate\_serial\_*} and
\texttt{ServingCostModel::estimate\_stage\_breakdown(...)} in
\texttt{astra-sim/workload/ServingCostModel.cc}.

\subsection{Serial Prefill and Decode}

Let \(B_p\) and \(C_p\) denote the configured prefill base latency and per-token
cost, and let \(s_c\) denote \texttt{compute\_scale}. The serial prefill model is:

\begin{equation*}
T_{\mathrm{prefill,serial}} =
s_c B_p + s_c N_p C_p
\end{equation*}

where \(N_p\) is the prompt-token count for the request.

Let \(B_d\) and \(C_d\) denote the configured decode base latency and per-token
cost, and let \(\phi_d\) denote
\texttt{cost\_model.decode.interference\_factor}. The serial decode-step model
for a decode batch of \(N_d\) tokens is:

\begin{equation*}
T_{\mathrm{decode,serial}} =
s_c \mathbf{1}_{\mathrm{first}} B_d + \phi_d s_c N_d C_d
\end{equation*}

In the serial runtime, \(N_d = 1\) per emitted token.

\subsection{Batched Prefill and Decode}

For colocated or PD batch scheduling, the legacy-compatible batched terms use
\texttt{batch\_efficiency} and decode \texttt{interference\_factor}. Let
\(\eta_p\) and \(\eta_d\) denote the prefill and decode batch-efficiency
scalars. With \(N_{\mathrm{batch}}\) equal to the sum of tokens in the batch:

\begin{equation*}
T_{\mathrm{prefill,batch}} =
s_c B_p + \eta_p s_c N_{\mathrm{batch}} C_p
\end{equation*}

\begin{equation*}
T_{\mathrm{decode,batch}} =
s_c \mathbf{1}_{\mathrm{first}} B_d +
\eta_d \phi_d s_c N_{\mathrm{batch}} C_d
\end{equation*}

When the component-wise stage fields are left unset, these equations are the
exact fallback used by \texttt{ServingCostModel::estimate\_stage\_breakdown(...)}.

\subsection{Compatibility-Mode KV Transfer}

For PD transfer, the default prompt-to-KV byte estimate comes from
\texttt{ServingCostModel::estimate\_kv\_transfer\_bytes(...)}. If no override is
provided, the simulator uses:

\begin{equation*}
Q_{\mathrm{KV}} =
N_p L H_{\mathrm{kv}} D_h \cdot 2 b_{\mathrm{kv}}
\end{equation*}

If \texttt{override\_bytes\_per\_prompt\_token} is enabled, the simulator uses:

\begin{equation*}
Q_{\mathrm{KV}} =
N_p q_{\mathrm{kv,override}}
\end{equation*}

The transfer delay itself is:

\begin{equation*}
T_{\mathrm{transfer}} =
L_{\mathrm{pd}} +
\frac{Q_{\mathrm{KV}}}{BW_{\mathrm{pd}} \epsilon_{\mathrm{pd}}}
\end{equation*}

where \(L_{\mathrm{pd}}\) is the configured transfer latency,
\(BW_{\mathrm{pd}}\) is the configured bandwidth, and
\(\epsilon_{\mathrm{pd}}\) is the configured efficiency.

\subsection{Optional Collective Addenda}

If \texttt{prefill\_collective} or \texttt{decode\_collective} is enabled, the
runtime issues a real ASTRA collective after the stage compute term. The
requested collective size is first scaled by \texttt{comm\_scale}:

\begin{equation*}
Q'_{\mathrm{collective}} =
s_m Q_{\mathrm{collective}}
\end{equation*}

The realized extra latency is then the collective completion time returned by
the ASTRA system model rather than a closed-form serving-side estimate.
